% Generated draft: requires mathematical and source review.
\documentclass[11pt]{article}
\usepackage[margin=25mm]{geometry}
\usepackage{fontspec}
\setmainfont{Latin Modern Roman}
\usepackage{amsmath,amssymb,mathtools,unicode-math}
\setmathfont{Latin Modern Math}
\usepackage{enumitem,needspace,xurl,hyperref}
\setlength{\parindent}{0pt}
\setlength{\parskip}{5pt}
\setlength{\emergencystretch}{3em}
\newcommand{\N}{\mathbb{N}}
\newcommand{\Z}{\mathbb{Z}}
\newcommand{\Q}{\mathbb{Q}}
\newcommand{\R}{\mathbb{R}}
\newcommand{\C}{\mathbb{C}}
\newcommand{\F}{\mathbb{F}}
\newcommand{\A}{\mathbb{A}}
\newcommand{\PP}{\mathbb{P}}
\newcommand{\E}{\mathbb{E}}
\newcommand{\eps}{\varepsilon}
\newcommand{\dd}{\,\mathrm{d}}
\begin{document}
\section*{UnsolvedMath ID: 30001523}
\textbf{Source problem code:} catalogue entry
\subsection{Definitions and mathematical statement}
Let \(p\) be a prime and let \(P\) be a finite \(p\)-group. Write \(Z(P)\) for its center and set
\[
C(P)=\{z\in Z(P):z^p=1\},\qquad c(P)=\dim_{\F_p}C(P).
\]
Here an elementary abelian \(p\)-group is an abelian group in which every nonidentity element has order \(p\), viewed as a vector space over \(\F_p\). Thus \(C(P)\) is the largest central elementary abelian \(p\)-subgroup of \(P\), and \(c(P)\) is its rank. Define
\[
n(P)=\min\{\dim_{\C}V: V\text{ is a finite-dimensional complex vector space and there exists an injective homomorphism }P\longrightarrow\operatorname{GL}(V)\}.
\]
The supplied question asks whether, for every prime \(p\) and every finite \(p\)-group \(P\),
\[
e(P)\leq 2n(P)-c(P).
\]
\textbf{Unresolved qualification.} The supplied record never defines \(e(P)\). Its title refers to cohomological exponents, but does not specify a cohomology theory, coefficients, grading, or the construction of this invariant. Consequently the displayed inequality records the supplied question but is not yet a fully specified mathematical proposition. No definition of \(e(P)\) is inferred here. The record states no restriction excluding the trivial group or any prime; the missing definition must also determine the meaning of \(e(P)\) at these boundary cases.
\subsection{Short English statement}
The supplied question asks whether an undefined invariant \(e(P)\) of every finite \(p\)-group is bounded above by twice its smallest faithful complex representation dimension minus the rank of its largest central elementary abelian subgroup. Defining \(e(P)\) is necessary before the question can be stated precisely. The supplied open status is a curation label, not an independently verified mathematical conclusion.
\subsection{Sources}
\begin{itemize}
\item Local UnsolvedMath registry, supplied record with permanent ID \(30001523\), source problem code \texttt{catalogue entry}, title ``Bounds on Cohomological Exponents of Finite \(p\)-Groups,'' in the collection ``Oberwolfach Reports Open Problems.'' No individual proposer is supplied; the proposed year is recorded as 2010.
\item Supplied source citation: \emph{Cohomology of Finite Groups: Interactions and Applications} (2011), Oberwolfach Reports, pp.~1885--1952, DOI \texttt{10.4171/owr/2010/32}. Supplied source URL: \url{https://doi.org/10.4171/owr/2010/32}. The background attributes the record to workshop open-problem discussions documented in Oberwolfach Reports.
\item Supplied literature reference: ``Oberwolfach Report 32, Cohomology of Finite Groups,'' \url{https://doi.org/10.4171/owr/2010/32}.
\item Literature assessment dated exactly \texttt{2026-08-21}: supplied dated curation, not independently verified. It labels the problem open under a stated classification rule because it says the extracted statement omits the conjectured inequality and retains only its setup. However, the supplied statement fields do contain the displayed inequality. This discrepancy remains unresolved. The background expressly calls for independent verification of time-sensitive status and cited claims.
\item The registry also supplies the statement-status label \texttt{corrected\_verified} and reports a manual semantic, typographic, and LaTeX review claiming that the cleanup is self-contained and adds no unsupported source-dependent content. That supplied review does not resolve the missing definition of \(e(P)\). The cited report has not been independently consulted for this draft.
\end{itemize}
\end{document}
